\documentclass[11pt,a4paper]{article}
\usepackage[T1]{fontenc}
\usepackage{lmodern}
\usepackage[margin=26mm]{geometry}
\usepackage{amsmath,amssymb,amsthm,booktabs,array,longtable}
\usepackage{microtype,multicol}
\usepackage[hidelinks]{hyperref}
\setlength{\parskip}{3pt}
\setlength{\emergencystretch}{2em}
\raggedbottom
\allowdisplaybreaks[1]
\clubpenalty=10000
\widowpenalty=10000
\newtheorem{theorem}{Theorem}[section]
\newtheorem{lemma}[theorem]{Lemma}
\newtheorem{corollary}[theorem]{Corollary}
\theoremstyle{remark}
\newtheorem{remark}[theorem]{Remark}
\newcommand{\Z}{\mathbb Z}
\newcommand{\Q}{\mathbb Q}
\newcommand{\N}{\mathbb N}
\newcommand{\F}{\mathbb F}
\newcommand{\ps}[1]{#1[[x]]}
\newcommand{\Rev}{\operatorname{Rev}}
\newcommand{\rev}{^{\langle-1\rangle}}
\newcommand{\Plus}{\operatorname{Plus}}
\newcommand{\Minus}{\operatorname{Minus}}
\newcommand{\coeff}[2]{[x^{#1}]#2}
\newcommand{\E}{\mathcal E}
\Urlmuskip=0mu plus 1mu
\hypersetup{pdftitle={A392210: Odd-index evenness in a nested binomial equation},pdfauthor={Dev Goyal},pdfsubject={Separate manuscript with AI disclosure and Lean proof map}}
\title{A392210\\[4pt]\large Odd-index evenness in a nested binomial equation}
\author{Dev Goyal}
\date{24 October 2026}
\begin{document}
\maketitle
\begin{center}
\fbox{\begin{minipage}{0.9\linewidth}\small
\textbf{AI disclosure.} ChatGPT (OpenAI) generated the mathematical arguments,
Lean proofs, and this manuscript in response to Dev Goyal's requests.
Dev Goyal is the sole human author. The source sequence and conjecture or
observation are credited to Paul D. Hanna.
\textbf{Independent human review is pending.}
\end{minipage}}
\end{center}
\begin{abstract}
The normalized equation $D=x+(D(x/(1-x)))^2$ has a unique integral solution. Every odd-index coefficient beyond the first is even. We prove the source construction and the elementary characteristic-two square argument; the result formalizes an observed pattern, not a new general principle.
\end{abstract}

\section{Source, notation, and scope}

The source is Paul D. Hanna's OEIS entry A392210 \cite{oeis210},
retrieved on 24 October 2026. Its initial sequence index is 1.
We also retain the formal zero constant coefficient at index zero.

All series are formal. The set $\N$ includes zero, and $[x^n]U$ denotes a
coefficient. A normalized series is $x+O(x^2)$; $O(x^r)$ means divisibility
by $x^r$, over any coefficient ring used below. We write $U\circ V=U(V(x))$
only when $V(0)=0$. Compositional inversion, denoted $\Rev(U)$ or $U\rev$,
is distinct from multiplicative inversion of a unit. Congruences are
coefficientwise; ordinary products and compositional iterates are distinct.

\section{Formal algebra used in the proof}

\begin{lemma}[Triangular construction]\label{lem:construction}
Let $K$ be a commutative ring. Suppose $\Phi$ preserves normalized series
and improves agreement by one degree:
\[
 G-H=O(x^r)\quad\Longrightarrow\quad
 \Phi(G)-\Phi(H)=O(x^{r+1})\qquad(r\ge2).
\]
Then $\Phi$ has a unique normalized fixed point. The same conclusion holds
for a family of series if the degree-$n$ equations, at every family index,
determine the degree-$n$ coefficients from strictly lower degrees.
\end{lemma}

\begin{proof}
The degree-$n$ coefficient of $\Phi(G)$ depends only on coefficients of
$G$ below degree $n$. Fix the constant and linear coefficients, and then
define every later coefficient successively. This gives a fixed point;
the same induction proves uniqueness. For a family, perform the induction
on degree simultaneously at every family index. The same induction applies with only the constant coefficient prescribed,
provided agreement is improved starting at degree one. It also applies to
an operator on all series when agreement is improved starting at degree zero.
For the operators used in this paper, the recurrence requires
only ring operations, so it also constructs integral solutions when
$K=\Z$. Equivalently, iterating $\Phi$ from any normalized initial series
fixes one further coefficient at each step.
\end{proof}

\begin{lemma}[Inversion and transport]\label{lem:inverse}
A normalized series over any commutative ring has a unique normalized
two-sided compositional inverse. A constant-one series has a unique
multiplicative inverse with constant coefficient one. These inverses and
valid composition commute with coefficient ring homomorphisms.
Normalized compositional inversion preserves agreement modulo $x^r$.
\end{lemma}

\begin{proof}
For normalized $G$, the degree-$n$ equation in $G(H)=x$ determines $h_n$
with multiplier one from earlier coefficients. Construct a right inverse
$J$ of $H$. Associativity gives
$G=G\circ(H\circ J)=(G\circ H)\circ J=J$, so $H\circ G=x$ as well.
The same associativity argument proves uniqueness. Multiplicative
inversion follows similarly from the coefficient recurrence for $UV=1$.

Each coefficient of a valid composition is a finite polynomial in the
coefficients of its arguments. Ring homomorphisms preserve these
polynomials and the defining equations, so uniqueness gives transport of
inversion. If $G-H=O(x^r)$, composition with any normalized series
preserves that agreement. Composing with the respective inverses then
shows $\Rev(G)-\Rev(H)=O(x^r)$.
\end{proof}

\section{A square-substitution construction}

\begin{lemma}[A useful square equation]\label{lem:square}
For normalized $T,P\in K[[x]]$, there is a unique normalized solution of
\[
 G=T+(G\circ P)^2.
\]
The construction commutes with coefficient ring homomorphisms.
\end{lemma}

\begin{proof}
If $G-H=O(x^r)$, then $(G-H)\circ P=O(x^r)$. The sum
$G\circ P+H\circ P$ has zero constant coefficient, so the difference
of the two squares is $O(x^{r+1})$. Normalization is preserved.
Lemma~\ref{lem:construction} gives existence and uniqueness; mapping the
equation and using uniqueness gives transport.
\end{proof}

\section{An elementary additional observation: A392210}\label{sec:odd}

The source gives the normalized functional equation \cite{oeis210}
\[
 D=x+\left(D\!\left(\frac{x}{1-x}\right)\right)^2. \tag{1}
\]
It describes evenness at odd indices beyond the first as an apparent
pattern. This follows directly from the square on the right.

\begin{corollary}
Equation (1) has a unique normalized integral solution. Every
odd-index coefficient beyond the first is even; equivalently,
$[x^{2n-1}]D$ is even for every $n>1$.
\end{corollary}

\begin{proof}
The normalized series $M=x/(1-x)$ is defined by the actual unit inverse
of $1-x$. Lemma~\ref{lem:square}, with $T=x$ and $P=M$, constructs $D$
and proves uniqueness for (1). Over $\F_2$, the square of any formal
series has only even exponents. Therefore the only possibly nonzero
odd-index coefficient on the right of (1) is the coefficient of $x$,
which is one. All later odd-index coefficients of $D$ vanish modulo two.
\end{proof}


\section{Separate Lean proof and checked scope}

The entry file is \nolinkurl{lean/Entries/A392210.lean}. Its explicit
theorem statements and proof terms use the namespace
\nolinkurl{FinalCollection.A392210}. The complete imported implementations
and their shared mathematical dependencies are included in the collection.
The following local theorem names give the proof map:

\begin{multicols}{2}\raggedright\begin{itemize}\setlength{\itemsep}{2pt}\setlength{\parsep}{0pt}
\item \nolinkurl{third_existence}.
\item \nolinkurl{third_definition}.
\item \nolinkurl{third_evenness}.
\item \nolinkurl{third_arbitrary_source}.
\end{itemize}\end{multicols}

The main implementation modules are
\begin{quote}\raggedright
\noindent\nolinkurl{Round4.NestedBinomial}\par
\noindent\nolinkurl{Round4.SquareComposition}\par
\end{quote}
The catalogue records the full original endpoint names and the new entry
theorems. The project pins Lean 4.19.0 and Mathlib revision
\nolinkurl{c44e0c8ee63ca166450922a373c7409c5d26b00b}. The collection
report records source recompilation and transitive axiom checks. Accepted
axioms are propositional extensionality, classical choice, and quotient
soundness. All index conclusions above are unbounded; exact finite
coefficient checks in the package provide separate corroboration.

\paragraph{Limits and status.}
This is the source apparent pattern and a direct characteristic-two square consequence. No analytic interpretation of the nested title is claimed.
The manuscript consolidates completed results; it does not assert global
novelty, exclusive priority, peer review, OEIS acceptance, or a new submission.
Standard algebraic ingredients are not claimed as new discoveries.

\begin{thebibliography}{9}
\raggedright
\bibitem{oeis210}
Paul D. Hanna, \emph{OEIS A392210}, 2 February 2026.
\url{https://oeis.org/A392210}. Retrieved 24 October 2026.
\end{thebibliography}
\end{document}
